% year: תשפ"ג
% type: מבחן
% semester: א
% moed: ב
% number: 2א
% points: 10
% categories: derivatives
% source: exams/מבחנים/תשפג/מועד ב תשפג.docx

\begin{question}
(10 נק') חשבו את משוואת המשיק למעגל $x^2+y^2=25$ בנקודה $(3,4)$.
\end{question}

\begin{hint}
גזרו את המשוואה בצורה סתומה.
\end{hint}

\begin{solution}
הנקודה על המעגל: $9+16=25$. בסביבת $(3,4)$ ($y>0$) המעגל הוא גרף של פונקציה גזירה $y(x)$. נגזור את המשוואה בצורה סתומה לפי $x$:
\[ 2x+2y\,y'=0\ \Longrightarrow\ y'=-\frac{x}{y},\qquad y'(3)=-\frac34. \]
משוואת המשיק: $y-4=-\frac34(x-3)$, כלומר
\[ y=-\frac34x+\frac{25}{4}\qquad\Longleftrightarrow\qquad 3x+4y=25. \]
(בהתאם לגאומטריה: המשיק ניצב לרדיוס בכיוון $(3,4)$.)
\end{solution}
